% SEGMENT OLP-0559-S01
\documentclass[../../../include/open-logic-section]{subfiles}

\begin{document}

\olfileid{sth}{spine}{valpha}

\olsection{படிநிலைகளை $V_\alpha$-களாக வரையறுத்தல்}

\olref[sth][ordinals][]{chap} இல், நன்கு வரிசைப்படுத்தல்களையும்
(ஃபான் நாய்மனின்) வரிசையெண்களையும் வரையறுத்தோம். இந்த இயலில், கணங்களின்
\emph{படிநிலைமுறையையே} வருணிக்க அவற்றைப் பயன்படுத்துவோம். இதைச் செய்ய,
\olref[sth][ordinals][opps]{sec} இல் பின்வரிசை வரிசையெண்கள், எல்லை
வரிசையெண்கள் என்ற கருத்துகளை வரையறுத்ததை நினைவுகூர்க. பின்வரும்
வரையறையில் இக்கருத்துகளைப் பயன்படுத்துகிறோம்:

\begin{defn}\ollabel{defValphas}
	\begin{align*}
	V_\emptyset &\defis \emptyset\\
	V_{\ordsucc{\alpha}} &\defis \Pow{V_\alpha} & & 
	\text{எந்த வரிசையெண் }\alpha\text{ க்கும்}\\
	V_{\alpha} &\defis \bigcup_{\gamma < \alpha} V_\gamma & & 
	\text{எப்போது }\alpha\text{ ஓர் எல்லை வரிசையெண்ணோ அப்போது}
\end{align*}
\end{defn}
\noindent
இது வரிசையெண்கள்மீது \emph{முடிவிலிக்கடந்த மீள்வரையறையால்} வழங்கப்படும்
வரையறை. இவ்வகையில், இயல் எண்கள்மீதான சார்புகளின் மீள்வரையறை
வரையறைகளுடன் இதை ஒப்பிட வேண்டும்.\footnote{\olref[sfr][infinite][dedekind]{sec}
இலுள்ள கூட்டல், பெருக்கல், அடுக்கேற்றம் ஆகியவற்றின் வரையறைகளுடன் ஒப்பிடுக.}
இயல் எண்களைக் கையாளும்போது போலவே, ஓர் அடிப்படை நேர்வையும் பின்வரிசை
நேர்வுகளையும் வரையறுக்கிறோம்; ஆனால் வரிசையெண்களைக் கையாளும்போது,
\emph{எல்லை} நேர்வுகளின் நடத்தையையும் விவரிக்க வேண்டும்.

$V_\alpha$-களின் இவ்வரையறை ஒரு முக்கியமான கட்டமாக இருக்கும். கணங்களின்
படிநிலைமுறையை, கணங்கள் \emph{படிநிலைகளாக} உருவாவதாக முறைசாராமல்
ஊக்குவித்தோம். நடைமுறையில், $V_\alpha$-களே அந்தப் படிநிலைகள்.
ஆனால் முக்கியமாக, இது படிநிலைகளின் \emph{அக} வருணனை. கோட்பாட்டின்
சாத்தியமான ஒரு \emph{மாதிரியை} முன்வைப்பதற்குப் பதிலாக, நமது
\emph{கணக் கோட்பாட்டிற்குள்ளேயே} படிநிலைகளை வரையறுத்திருப்போம்.

\end{document}
